\section*{Chapter \ref{ch:g12:continental-crust} --- The Continental Crust}

\begin{solution}{exo:g12:continental-crust:1}
\omterm{def:g12:continental-crust:continental-crust}{Continental crust}: 30 to \qty{70}{km}, about \qty{41}{km} on average;
granitic rocks, gneisses and other metamorphic and \omterm{def:g11:igneous-rocks:volcanic-plutonic}{plutonic rocks} under a
sedimentary cover; about \qty{61}{\percent} silica on average; about
\qty{2830}{kg/m^3}; up to 4.03 billion years (grains 4.4). \omterm{def:g11:ocean-ridges:oceanic-lithosphere}{Oceanic crust}:
about \qty{7}{km}; basalt and \omterm{def:g11:ocean-ridges:gabbro}{gabbro}; about \qty{50}{\percent} silica; about
\qty{3000}{kg/m^3}; less than 200 million years.
\end{solution}

\begin{solution}{exo:g12:continental-crust:2}
\qty{60.6}{\percent}: \omterm{def:g11:igneous-rocks:felsic-mafic}{intermediate}, like a \omterm{def:g11:igneous-rocks:rock-names}{diorite}. \qty{66.6}{\percent}:
\omterm{def:g11:igneous-rocks:felsic-mafic}{felsic}, like a granite.
\end{solution}

\begin{solution}{exo:g12:continental-crust:3}
The \omterm{def:g11:ocean-ridges:oceanic-lithosphere}{oceanic crust}, about \qty{3000}{kg/m^3}, sits on lithospheric mantle,
and the old, cold \omterm{def:g11:ocean-ridges:oceanic-lithosphere}{oceanic lithosphere} as a whole is dense enough to sink.
The \omterm{def:g12:continental-crust:continental-crust}{continental crust}, about \qty{2830}{kg/m^3}, far lighter than the mantle
(about \qty{3380}{kg/m^3}), is too buoyant: a continent arriving at a trench
stops \omterm{def:g11:subduction-zones:subduction}{subduction}.
\end{solution}

\begin{solution}{exo:g12:continental-crust:4}
(1) An old, stable core of a continent, which no longer deforms. (2) More
than a billion years. (3) Metamorphic and \omterm{def:g11:igneous-rocks:volcanic-plutonic}{plutonic rocks}, the deep \omterm{def:g12:continental-crust:root}{roots} of
former mountains.
\end{solution}

\begin{solution}{exo:g12:continental-crust:5}
$4030/190 \approx 21$ times. The ocean floor is made at the ridges and
destroyed at \omterm{def:g11:subduction-zones:subduction}{subduction zones}: none survives much more than 200 million
years.
\end{solution}

\begin{solution}{exo:g12:continental-crust:6}
(1) $0.306 \times 210.4 \approx \qty{64.4}{million~km^2}$.
(2) $210.4 - 64.4 = \qty{146.0}{million~km^2}$.
(3) $4\pi \times 6371^2 \approx \qty{5.10e8}{km^2}$; $210.4/510.1 \approx
0.41$: \qty{41}{\percent}.
\end{solution}

\begin{solution}{exo:g12:continental-crust:7}
(1) $41100 \times 2830 \approx \qty{1.16e8}{kg}$. (2) $7100 \times 3000 =
\qty{2.13e7}{kg}$. (3) About 5.5 times.
\end{solution}

\begin{solution}{exo:g12:continental-crust:8}
\omterm{def:g12:continental-crust:root}{Root}: $2.5 \times 2830/550 \approx \qty{12.9}{km}$. The crust is thicker by
$2.5 \times 3380/550 \approx \qty{15.4}{km}$: about \qty{50}{km} in all,
within the 50 to \qty{70}{km} found under mountain belts.
\end{solution}

\begin{solution}{exo:g12:continental-crust:9}
(1) About $0.103 + 0.100 + 0.102 \approx 0.30$: \qty{30}{\percent}.
(2) The youngest band, 0 to 10 million years (about \qty{10}{\percent}),
almost tied with 20 to 30. (3) Old floor has been destroyed at \omterm{def:g11:subduction-zones:subduction}{subduction
zones}; the older the floor, the more of it has gone.
\end{solution}

\begin{solution}{exo:g12:continental-crust:10}
$120 = 0.694\,h - 0.306 \times 100$, so $h = 150.6/0.694 \approx
\qty{217}{m}$. Counting only the land leaves out the low, submerged edges of
the continents, and the mean comes out far higher.
\end{solution}

\begin{solution}{exo:g12:continental-crust:11}
The island: a \omterm{def:g11:ocean-ridges:gabbro}{gabbro} under lavas and a \omterm{def:g11:seismic-structure:moho}{Moho} at \qty{8}{km}, \omterm{def:g11:ocean-ridges:oceanic-lithosphere}{oceanic crust}
(the elevation of the island does not change that). The plain: gneiss and a
\omterm{def:g11:seismic-structure:moho}{Moho} at \qty{38}{km}, \omterm{def:g12:continental-crust:continental-crust}{continental crust}.
\end{solution}

\begin{solution}{exo:g12:continental-crust:12}
(1) Exact: $\tfrac{4}{3}\pi(6371^3 - 6329.9^3) \approx
\qty{2.0829e10}{km^3}$; approximate: $4\pi \times 6371^2 \times 41.1
\approx \qty{2.0964e10}{km^3}$. (2) The approximation is too large by about
\qty{0.65}{\percent}, close to $H/R \approx 0.0065$: for a thin layer the
error is negligible.
\end{solution}

\begin{solution}{exo:g12:continental-crust:13}
(1) $60.6 = (66.6 + 2x)/3$, so $x = (181.8 - 66.6)/2 = \qty{57.6}{\percent}$.
(2) \omterm{def:g11:igneous-rocks:felsic-mafic}{Intermediate}. (3) The deep \omterm{def:g12:continental-crust:continental-crust}{continental crust} is less rich in silica than
the upper crust: more \omterm{def:g11:igneous-rocks:felsic-mafic}{mafic} minerals, richer in iron, magnesium and calcium.
\end{solution}

\begin{solution}{exo:g12:continental-crust:14}
(1) $3 \times 10^8/3 = 10^8$ years, 100 million years. (2) The same order as
the ages of the ocean floor: three quarters younger than 100 million years,
none older than about 190. (3) It must be destroyed as fast as it is made:
it sinks back into the mantle at \omterm{def:g11:subduction-zones:subduction}{subduction zones}.
\end{solution}

\begin{solution}{exo:g12:continental-crust:15}
(1) $8.6 \times 10^9 / 4.0 \times 10^9 \approx \qty{2.2}{km^3}$ a year.
(2) A steady rate is only one possibility: the crust may have grown fast early
and been recycled since, or grown late. The oldest rocks and grains prove
that some crust existed early; they cannot tell how much, because old crust
can be eroded or carried back into the mantle and leave few traces.
\end{solution}

\begin{solution}{pb:g12:continental-crust:1}
\textbf{1.} $4\pi \times 6371^2 \approx \qty{5.10e8}{km^2}$.

\textbf{2.} $210.4/510.1 \approx 0.41$: \qty{41}{\percent}.

\textbf{3.} $0.306 \times 210.4 \approx \qty{64.4}{million~km^2}$.

\textbf{4.} $510.1 - 210.4 \approx \qty{299.7}{million~km^2}$,
\qty{59}{\percent}.

\textbf{5.} The edges of the continents lie under shallow seas, on the
continental shelves.

\textbf{6.} The thickness, \qty{41}{km}, is very small compared with the
radius, \qty{6371}{km}: the curvature changes the volume by less than one
percent.

\textbf{7.} $210.4 \times 10^6 \times 41.1 \approx \qty{8.65e9}{km^3}$.

\textbf{8.} \qty{8.65e18}{m^3}.

\textbf{9.} $210.4 \times 10^6 \times 36 \approx \qty{7.57e9}{km^3}$.

\textbf{10.} $(8.65 - 7.57)/7.57 \approx 0.14$: \qty{14}{\percent}.

\textbf{11.} $8.65 \times 10^{18} \times 2830 \approx \qty{2.45e22}{kg}$.

\textbf{12.} $41100 \times 2830 \approx \qty{1.16e8}{kg}$.

\textbf{13.} $299.7 \times 10^6 \times 7.1 \approx \qty{2.13e9}{km^3}$, that
is \qty{2.13e18}{m^3}; times 3000: about \qty{6.4e21}{kg}.

\textbf{14.} The \omterm{def:g12:continental-crust:continental-crust}{continental crust}, about 3.8 times heavier: it is almost six
times thicker, which outweighs its smaller area and density.

\textbf{15.} It is made of granitic rocks, rich in quartz and feldspars, the
lightest rock-forming minerals; the \omterm{def:g11:ocean-ridges:oceanic-lithosphere}{oceanic crust} is made of basalt and
\omterm{def:g11:ocean-ridges:gabbro}{gabbro}, rich in pyroxene and olivine.

\textbf{16.} $2.45 \times 10^{22}/5.97 \times 10^{24} \approx 0.0041$.

\textbf{17.} $3 \times 0.41 \times (41.1/6371) \times (2830/5513) \approx
0.0041$: the same.

\textbf{18.} $6.4 \times 10^{21}/5.97 \times 10^{24} \approx 0.0011$: about a
tenth of a percent.

\textbf{19.} About \qty{0.4}{\percent}.

\textbf{20.} The \omterm{def:g12:continental-crust:continental-crust}{continental crust} is about \textbf{\qty{0.4}{\percent} of the
Earth's mass}.
\end{solution}
